\documentclass[10pt,a4paper]{amsart}
\usepackage{amsmath,amssymb,amsthm,graphicx}
\usepackage{verbatim}
\newenvironment{reference}{\comment}{\endcomment}
%\newenvironment{reference}{}{}
\newenvironment{slogan}{\comment}{\endcomment}
\newenvironment{history}{\comment}{\endcomment}
% For commutative diagrams we use Xy-pic
\usepackage[all]{xy}

% We use 2cell for 2-commutative diagrams.
\xyoption{2cell}
% Theorem environments.
%
\theoremstyle{plain}
\newtheorem{theorem}[subsection]{Theorem}
\newtheorem{proposition}[subsection]{Proposition}
\newtheorem{lemma}[subsection]{Lemma}

\theoremstyle{definition}
\newtheorem{definition}[subsection]{Definition}
\newtheorem{example}[subsection]{Example}
\newtheorem{exercise}[subsection]{Exercise}
\newtheorem{situation}[subsection]{Situation}

\theoremstyle{remark}
\newtheorem{remark}[subsection]{Remark}
\newtheorem{remarks}[subsection]{Remarks}

\numberwithin{equation}{subsection}

% Macros
%
\def\lim{\mathop{\mathrm{lim}}\nolimits}
\def\colim{\mathop{\mathrm{colim}}\nolimits}
\def\Spec{\mathop{\mathrm{Spec}}}
\def\Hom{\mathop{\mathrm{Hom}}\nolimits}
\def\Ext{\mathop{\mathrm{Ext}}\nolimits}
\def\SheafHom{\mathop{\mathcal{H}\!\mathit{om}}\nolimits}
\def\SheafExt{\mathop{\mathcal{E}\!\mathit{xt}}\nolimits}
\def\Sch{\mathit{Sch}}
\def\Mor{\mathop{\mathrm{Mor}}\nolimits}
\def\Ob{\mathop{\mathrm{Ob}}\nolimits}
\def\Sh{\mathop{\mathit{Sh}}\nolimits}
\def\NL{\mathop{N\!L}\nolimits}
\def\CH{\mathop{\mathrm{CH}}\nolimits}
\def\proetale{{pro\text{-}\acute{e}tale}}
\def\etale{{\acute{e}tale}}
\def\QCoh{\mathit{QCoh}}
\def\Ker{\mathop{\mathrm{Ker}}}
\def\Im{\mathop{\mathrm{Im}}}
\def\Coker{\mathop{\mathrm{Coker}}}
\def\Coim{\mathop{\mathrm{Coim}}}

% Boxtimes
%
\DeclareMathSymbol{\boxtimes}{\mathbin}{AMSa}{"02}

%
% Macros for moduli stacks/spaces
%
\def\QCohstack{\mathcal{QC}\!\mathit{oh}}
\def\Cohstack{\mathcal{C}\!\mathit{oh}}
\def\Spacesstack{\mathcal{S}\!\mathit{paces}}
\def\Quotfunctor{\mathrm{Quot}}
\def\Hilbfunctor{\mathrm{Hilb}}
\def\Curvesstack{\mathcal{C}\!\mathit{urves}}
\def\Polarizedstack{\mathcal{P}\!\mathit{olarized}}
\def\Complexesstack{\mathcal{C}\!\mathit{omplexes}}
% \Pic is the operator that assigns to X its picard group, usage \Pic(X)
% \Picardstack_{X/B} denotes the Picard stack of X over B
% \Picardfunctor_{X/B} denotes the Picard functor of X over B
\def\Pic{\mathop{\mathrm{Pic}}\nolimits}
\def\Picardstack{\mathcal{P}\!\mathit{ic}}
\def\Picardfunctor{\mathrm{Pic}}
\def\Deformationcategory{\mathcal{D}\!\mathit{ef}}
\usepackage{cleveref}
\crefname{originallemma}{original lemma}{original lemmas}
\Crefname{originallemma}{Original lemma}{Original lemmas}
\makeatletter
\newlabel{lemma-base-change-gerbe}{{lemma-base-change-gerbe}{source}{}{sourcepointer-1}{}}
\newlabel{lemma-base-change-gerbe@cref}{{[originallemma][0][]lemma-base-change-gerbe}{[1][0][]source}{}{}{}}
\makeatother
\crefname{originallemma}{original lemma}{original lemmas}
\Crefname{originallemma}{Original lemma}{Original lemmas}
\makeatletter
\newlabel{lemma-cartesian-square-inertia}{{lemma-cartesian-square-inertia}{source}{}{sourcepointer-2}{}}
\newlabel{lemma-cartesian-square-inertia@cref}{{[originallemma][0][]lemma-cartesian-square-inertia}{[1][0][]source}{}{}{}}
\makeatother
\crefname{originallemma}{original lemma}{original lemmas}
\Crefname{originallemma}{Original lemma}{Original lemmas}
\makeatletter
\newlabel{lemma-descent-finite-presentation}{{lemma-descent-finite-presentation}{source}{}{sourcepointer-3}{}}
\newlabel{lemma-descent-finite-presentation@cref}{{[originallemma][0][]lemma-descent-finite-presentation}{[1][0][]source}{}{}{}}
\makeatother
\crefname{originallemma}{original lemma}{original lemmas}
\Crefname{originallemma}{Original lemma}{Original lemmas}
\makeatletter
\newlabel{lemma-descent-flat}{{lemma-descent-flat}{source}{}{sourcepointer-4}{}}
\newlabel{lemma-descent-flat@cref}{{[originallemma][0][]lemma-descent-flat}{[1][0][]source}{}{}{}}
\makeatother
\crefname{originallemma}{original lemma}{original lemmas}
\Crefname{originallemma}{Original lemma}{Original lemmas}
\makeatletter
\newlabel{lemma-local-structure-gerbe}{{lemma-local-structure-gerbe}{source}{}{sourcepointer-5}{}}
\newlabel{lemma-local-structure-gerbe@cref}{{[originallemma][0][]lemma-local-structure-gerbe}{[1][0][]source}{}{}{}}
\makeatother
\crefname{originalproposition}{original proposition}{original propositions}
\Crefname{originalproposition}{Original proposition}{Original propositions}
\makeatletter
\newlabel{proposition-when-gerbe}{{proposition-when-gerbe}{source}{}{sourcepointer-6}{}}
\newlabel{proposition-when-gerbe@cref}{{[originalproposition][0][]proposition-when-gerbe}{[1][0][]source}{}{}{}}
\makeatother
\title{Flatness and finite presentation of the inertia, and surjectivity of the diagonal, for a gerbe}
\author{The Stacks Project Authors}
\date{stacks-project:stacks-morphisms.tex:local-snapshot-20261012; GNU FDL 1.2 or later}
\begin{document}
\maketitle
\noindent\textit{Editorial note.} All delivered mathematical text is the author's own, copied from the pinned Stacks Project chapter source sources/stacks-algebraic-stacks-morphisms/pinned-original-source.tex; no statement or proof step was rewritten, reordered or removed.\par\medskip
\noindent\textit{Editorial note.} Preamble provenance: the Stacks Project keeps one single shared preamble for every chapter; this card takes it directly from the staged sources/stacks-algebraic-stacks-morphisms/preamble.tex as three verbatim spans (lines 8-12, 14-18 and 154-217), so no derived file was written and nothing was added or redefined. The wrapper class and the base packages are supplied by the assembler.\par\medskip
\noindent\textit{Editorial note.} Bibliography provenance: neither the counted statement nor its proof contains a citation (this lane rescans the whole span before writing the card), so the card carries no bibliography and no entry was reconstructed.\par\medskip
\noindent\textit{Editorial note.} Selection provenance: this card was opened under the file-based control contract (CONTROL.md), with the mechanical pre-card gate run on its statement and proof interval before the card was written.\par\medskip
\noindent\textit{Allowed standard prerequisites (editorial).}\begin{itemize}
\item The fibre product description of the inertia of a base change. Invoked inside the counted proof (source lines 6456-6459). Allows the inertia of the base-changed stack to be computed from that of the original.
\item The descent statements for flatness and for local finite presentation. Invoked inside the counted proof (source lines 6460-6462 and 6481-6484). Reduce the properties of the inertia and of the diagonal to a flat finitely presented cover.
\item The characterization of a gerbe by flatness and local finite presentation of its inertia. Invoked inside the counted proof (source lines 6463-6465). Gives flatness and local finite presentation of the inertia of the base-changed stack.
\item The local structure of a gerbe as a quotient stack by a flat group space. Invoked inside the counted proof (source lines 6470-6473). Realises the base-changed stack as a quotient in order to compute the diagonal.
\end{itemize}\medskip
\ifdefined\hypertarget\hypertarget{sourcepointer-1}{}\fi
\noindent\textit{Source reference.} Original lemma lemma-base-change-gerbe: sources/stacks-algebraic-stacks-morphisms/pinned-original-source.tex, source lines 6129--6129. Navigation pointer only: the counted text refers to the Stacks item lemma-base-change-gerbe, which is an existing labeled result of this same chapter and is not part of this card. The pointer carries the target's own label as its number and its definition line as its span; its content is not reproduced here.\par\medskip
\ifdefined\hypertarget\hypertarget{sourcepointer-2}{}\fi
\noindent\textit{Source reference.} Original lemma lemma-cartesian-square-inertia: sources/stacks-algebraic-stacks-morphisms/pinned-original-source.tex, source lines 1436--1436. Navigation pointer only: the counted text refers to the Stacks item lemma-cartesian-square-inertia, which is an existing labeled result of this same chapter and is not part of this card. The pointer carries the target's own label as its number and its definition line as its span; its content is not reproduced here.\par\medskip
\ifdefined\hypertarget\hypertarget{sourcepointer-3}{}\fi
\noindent\textit{Source reference.} Original lemma lemma-descent-finite-presentation: sources/stacks-algebraic-stacks-morphisms/pinned-original-source.tex, source lines 5833--5833. Navigation pointer only: the counted text refers to the Stacks item lemma-descent-finite-presentation, which is an existing labeled result of this same chapter and is not part of this card. The pointer carries the target's own label as its number and its definition line as its span; its content is not reproduced here.\par\medskip
\ifdefined\hypertarget\hypertarget{sourcepointer-4}{}\fi
\noindent\textit{Source reference.} Original lemma lemma-descent-flat: sources/stacks-algebraic-stacks-morphisms/pinned-original-source.tex, source lines 5389--5389. Navigation pointer only: the counted text refers to the Stacks item lemma-descent-flat, which is an existing labeled result of this same chapter and is not part of this card. The pointer carries the target's own label as its number and its definition line as its span; its content is not reproduced here.\par\medskip
\ifdefined\hypertarget\hypertarget{sourcepointer-5}{}\fi
\noindent\textit{Source reference.} Original lemma lemma-local-structure-gerbe: sources/stacks-algebraic-stacks-morphisms/pinned-original-source.tex, source lines 6248--6248. Navigation pointer only: the counted text refers to the Stacks item lemma-local-structure-gerbe, which is an existing labeled result of this same chapter and is not part of this card. The pointer carries the target's own label as its number and its definition line as its span; its content is not reproduced here.\par\medskip
\ifdefined\hypertarget\hypertarget{sourcepointer-6}{}\fi
\noindent\textit{Source reference.} Original proposition proposition-when-gerbe: sources/stacks-algebraic-stacks-morphisms/pinned-original-source.tex, source lines 6319--6319. Navigation pointer only: the counted text refers to the Stacks item proposition-when-gerbe, which is an existing labeled result of this same chapter and is not part of this card. The pointer carries the target's own label as its number and its definition line as its span; its content is not reproduced here.\par\medskip
\begin{lemma}
\label{lemma-gerbe-isom-fppf}
Let $f : \mathcal{X} \to \mathcal{Y}$ be a morphism of algebraic stacks
which makes $\mathcal{X}$ a gerbe over $\mathcal{Y}$. Then
\begin{enumerate}
\item $\mathcal{I}_{\mathcal{X}/\mathcal{Y}} \to \mathcal{X}$
is flat and locally of finite presentation,
\item $\mathcal{X} \to \mathcal{X} \times_\mathcal{Y} \mathcal{X}$
is surjective, flat, and locally of finite presentation,
\item given algebraic spaces $T_i$, $i = 1, 2$ and morphisms
$x_i : T_i \to \mathcal{X}$, with $y_i = f \circ x_i$ the morphism
$$
T_1 \times_{x_1, \mathcal{X}, x_2} T_2 \longrightarrow
T_1 \times_{y_1, \mathcal{Y}, y_2} T_2
$$
is surjective, flat, and locally of finite presentation,
\item given an algebraic space $T$ and morphisms
$x_i : T \to \mathcal{X}$, $i = 1, 2$, with $y_i = f \circ x_i$ the morphism
$$
\mathit{Isom}_\mathcal{X}(x_1, x_2) \longrightarrow
\mathit{Isom}_\mathcal{Y}(y_1, y_2)
$$
is surjective, flat, and locally of finite presentation.
\end{enumerate}
\end{lemma}

\begin{proof}
Proof of (1).
Choose a scheme $Y$ and a surjective smooth morphism $Y \to \mathcal{Y}$.
Set $\mathcal{X}' = \mathcal{X} \times_\mathcal{Y} Y$.
By Lemma \ref{lemma-cartesian-square-inertia} we obtain cartesian
squares
$$
\xymatrix{
\mathcal{I}_{\mathcal{X}'} \ar[r] \ar[d] &
\mathcal{X}' \ar[r] \ar[d] & Y \ar[d] \\
\mathcal{I}_{\mathcal{X}/\mathcal{Y}} \ar[r] &
\mathcal{X} \ar[r] & \mathcal{Y}
}
$$
By Lemmas \ref{lemma-descent-flat} and
\ref{lemma-descent-finite-presentation}
it suffices to prove that $\mathcal{I}_{\mathcal{X}'} \to \mathcal{X}'$
is flat and locally of finite presentation.
This follows from Proposition \ref{proposition-when-gerbe}
(because $\mathcal{X}'$ is a gerbe over $Y$ by
Lemma \ref{lemma-base-change-gerbe}).

\medskip\noindent
Proof of (2). With notation as above, note that we may assume that
$\mathcal{X}' = [Y/G]$ for some group algebraic space $G$ flat and
locally of finite presentation over $Y$, see
Lemma \ref{lemma-local-structure-gerbe}.
The base change of the morphism
$\Delta : \mathcal{X} \to \mathcal{X} \times_\mathcal{Y} \mathcal{X}$
over $\mathcal{Y}$ by the morphism $Y \to \mathcal{Y}$
is the morphism
$\Delta' : \mathcal{X}' \to \mathcal{X}' \times_Y \mathcal{X}'$.
Hence it suffices to show that $\Delta'$ is
surjective, flat, and locally of finite presentation
(see Lemmas \ref{lemma-descent-flat} and
\ref{lemma-descent-finite-presentation}).
In other words, we have to show that
$$
[Y/G] \longrightarrow [Y/G \times_Y G]
$$
is surjective, flat, and locally of finite presentation.
This is true because the base change by the surjective, flat,
locally finitely presented morphism $Y \to [Y/G \times_Y G]$
is the morphism $G \to Y$.

\medskip\noindent
Proof of (3). Observe that the diagram
$$
\xymatrix{
T_1 \times_{x_1, \mathcal{X}, x_2} T_2 \ar[d] \ar[r] &
T_1 \times_{y_1, \mathcal{Y}, y_2} T_2 \ar[d] \\
\mathcal{X} \ar[r] & \mathcal{X} \times_\mathcal{Y} \mathcal{X}
}
$$
is cartesian. Hence (3) follows from (2).

\medskip\noindent
Proof of (4). This is true because
$$
\mathit{Isom}_\mathcal{X}(x_1, x_2) =
(T \times_{x_1, \mathcal{X}, x_2} T) \times_{T \times T, \Delta_T} T
$$
hence the morphism in (4) is a base change of the morphism in (3).
\end{proof}

\end{document}
